\documentclass[tikz,border=8pt]{standalone}
\usepackage{tikz}
\usepackage{amsmath}
\usepackage{amssymb}
\usepackage{pgfplots}
\pgfplotsset{compat=1.18}
\usepackage{ctex}
\usetikzlibrary{calc,positioning,arrows.meta,matrix}

% LC 146 LRU 缓存设计
% 哈希表 + 双向链表：head ↔ A ↔ B ↔ C ↔ tail
% 三个操作快照：get / put 新 / put 满淘汰

\begin{document}
\begin{tikzpicture}[
  every node/.style={font=\small},
  dllnode/.style={draw=gray!70, rectangle, rounded corners=3pt,
    minimum width=1.1cm, minimum height=0.7cm, thick, fill=blue!8},
  headtail/.style={draw=gray!50, rectangle, rounded corners=3pt,
    minimum width=0.9cm, minimum height=0.7cm, thick, fill=gray!15},
  hashbox/.style={draw=orange!60, rectangle, rounded corners=3pt,
    minimum width=2.4cm, minimum height=0.7cm, thick, fill=orange!8},
  highlight/.style={draw=red!60!black, rectangle, rounded corners=3pt,
    minimum width=1.1cm, minimum height=0.7cm, very thick, fill=red!10},
  greennode/.style={draw=green!60!black, rectangle, rounded corners=3pt,
    minimum width=1.1cm, minimum height=0.7cm, very thick, fill=green!10},
  ptr/.style={->,>=Stealth,thick},
  biarr/.style={<->,>=Stealth,thick,gray!60},
]

%% ── 标题 ──────────────────────────────────────────────────
\node[font=\large\bfseries] at (6.5, 1.0)
  {LRU 缓存（LC 146 LRU Cache）};
\node[font=\small,gray] at (6.5, 0.3)
  {哈希表 + 双向链表；最近使用移到链表头，淘汰从链表尾移除};

%% ══════════════════════════════════════════════════════════
%% 快照 A：初始状态 capacity=3，已存 A(1) B(2) C(3)
%% 链表：head ↔ C ↔ B ↔ A ↔ tail（C 最新，A 最旧）
%% ══════════════════════════════════════════════════════════
\def\yA{-0.5}
\node[font=\small\bfseries,blue!70] at (0.2, \yA) {快照 1：初始（容量=3）};

% 双向链表
\node[headtail] (hA) at (3.0, \yA-0.8) {head};
\node[dllnode]  (cA) at (4.8, \yA-0.8) {C:3};
\node[dllnode]  (bA) at (6.4, \yA-0.8) {B:2};
\node[dllnode]  (aA) at (8.0, \yA-0.8) {A:1};
\node[headtail] (tA) at (9.8, \yA-0.8) {tail};

\draw[biarr] (hA.east) -- (cA.west);
\draw[biarr] (cA.east) -- (bA.west);
\draw[biarr] (bA.east) -- (aA.west);
\draw[biarr] (aA.east) -- (tA.west);

% 说明
\node[font=\tiny,gray,anchor=north] at (cA.south) {最近};
\node[font=\tiny,gray,anchor=north] at (aA.south) {最旧};

% 哈希表
\node[font=\scriptsize,anchor=east] at (2.8, \yA-2.0) {HashMap:};
\node[hashbox] (hm1A) at (4.0, \yA-2.0) {key=1 → A};
\node[hashbox] (hm2A) at (6.5, \yA-2.0) {key=2 → B};
\node[hashbox] (hm3A) at (9.0, \yA-2.0) {key=3 → C};

% 哈希 → 节点指针（斜线）
\draw[->,>=Stealth,orange!60,dashed,thin] (hm1A.north) -- (aA.south);
\draw[->,>=Stealth,orange!60,dashed,thin] (hm2A.north) -- (bA.south);
\draw[->,>=Stealth,orange!60,dashed,thin] (hm3A.north) -- (cA.south);

%% ══════════════════════════════════════════════════════════
%% 快照 B：get(1) → 将 A 移到头部，A 成最新
%% 链表：head ↔ A ↔ C ↔ B ↔ tail
%% ══════════════════════════════════════════════════════════
\def\yB{-4.0}
\node[font=\small\bfseries,blue!70] at (0.2, \yB) {快照 2：get(1)→移 A 到头};

\node[headtail] (hB) at (3.0, \yB-0.8) {head};
\node[greennode] (aB) at (4.8, \yB-0.8) {A:1};
\node[dllnode]  (cB) at (6.4, \yB-0.8) {C:3};
\node[dllnode]  (bB) at (8.0, \yB-0.8) {B:2};
\node[headtail] (tB) at (9.8, \yB-0.8) {tail};

\draw[biarr] (hB.east) -- (aB.west);
\draw[biarr] (aB.east) -- (cB.west);
\draw[biarr] (cB.east) -- (bB.west);
\draw[biarr] (bB.east) -- (tB.west);

\node[font=\tiny,gray,anchor=north] at (aB.south) {最近};
\node[font=\tiny,gray,anchor=north] at (bB.south) {最旧};

% 哈希表不变
\node[font=\scriptsize,anchor=east] at (2.8, \yB-2.0) {HashMap:};
\node[hashbox] (hm1B) at (4.0, \yB-2.0) {key=1 → A};
\node[hashbox] (hm2B) at (6.5, \yB-2.0) {key=2 → B};
\node[hashbox] (hm3B) at (9.0, \yB-2.0) {key=3 → C};

\draw[->,>=Stealth,orange!60,dashed,thin] (hm1B.north) -- (aB.south);
\draw[->,>=Stealth,orange!60,dashed,thin] (hm2B.north) -- (bB.south);
\draw[->,>=Stealth,orange!60,dashed,thin] (hm3B.north) -- (cB.south);

% 操作说明
\node[draw=gray!40,fill=white,thin,rounded corners=2pt,
      font=\scriptsize,inner sep=4pt,align=left]
  at (12.5, \yB-0.8)
  {get(1)：\\1. HashMap 查 key=1\\2. 找到节点 A\\3. 从原位置摘除\\4. 插到 head 后};

%% ══════════════════════════════════════════════════════════
%% 快照 C：put(4,v) 满了 → 淘汰 B（tail 前一个），插入 D
%% 链表：head ↔ D ↔ A ↔ C ↔ tail
%% ══════════════════════════════════════════════════════════
\def\yC{-7.5}
\node[font=\small\bfseries,blue!70] at (0.2, \yC) {快照 3：put(4,v)→满，淘汰 B，插入 D};

\node[headtail] (hC) at (3.0, \yC-0.8) {head};
\node[greennode] (dC) at (4.8, \yC-0.8) {D:4};
\node[dllnode]  (aC) at (6.4, \yC-0.8) {A:1};
\node[dllnode]  (cC) at (8.0, \yC-0.8) {C:3};
\node[headtail] (tC) at (9.8, \yC-0.8) {tail};

\draw[biarr] (hC.east) -- (dC.west);
\draw[biarr] (dC.east) -- (aC.west);
\draw[biarr] (aC.east) -- (cC.west);
\draw[biarr] (cC.east) -- (tC.west);

\node[font=\tiny,gray,anchor=north] at (dC.south) {最近};
\node[font=\tiny,gray,anchor=north] at (cC.south) {最旧};

% 哈希表：删 key=2，加 key=4
\node[font=\scriptsize,anchor=east] at (2.8, \yC-2.0) {HashMap:};
\node[hashbox] (hm1C) at (4.0, \yC-2.0) {key=1 → A};
\node[hashbox,draw=red!50,fill=red!5] (hm2C) at (6.5, \yC-2.0) {key=2 删除};
\node[hashbox] (hm3C) at (9.0, \yC-2.0) {key=3 → C};
\node[hashbox,draw=green!60!black,fill=green!8] (hm4C) at (11.5, \yC-2.0) {key=4 → D};

\draw[->,>=Stealth,orange!60,dashed,thin] (hm1C.north) -- (aC.south);
\draw[->,>=Stealth,orange!60,dashed,thin] (hm3C.north) -- (cC.south);
\draw[->,>=Stealth,green!60!black,dashed,thin] (hm4C.north) to[bend left=20] (dC.south);

% 淘汰说明
\node[draw=gray!40,fill=white,thin,rounded corners=2pt,
      font=\scriptsize,inner sep=4pt,align=left]
  at (12.5, \yC-0.8)
  {put(4,v)：\\1. size==capacity\\2. 淘汰 tail 前节点 B\\3. HashMap 删 key=2\\4. 新建节点 D\\5. 插到 head 后};

%% ── 分隔线 ──────────────────────────────────────────────
\draw[gray!25,dashed] (0.0, -10.8) -- (14.0, -10.8);

%% ── 算法说明 ─────────────────────────────────────────────
\node[draw=gray!50,fill=white,rounded corners=3pt,font=\small,
      inner sep=8pt,align=left]
  at (7.0, -11.8)
  {\textbf{数据结构：}HashMap$\langle$key, Node$\rangle$ + 带哨兵头尾的双向链表\\[4pt]
   \textbf{get(key)：}O(1) 查表 $\to$ 节点移到 head 后 $\to$ 返回 val（未找到返回 -1）\\[2pt]
   \textbf{put(key, val)：}O(1) 若 key 存在则更新并移头；否则新建节点插头；\\
   \quad\quad\quad\quad\quad\quad\quad 若超容量则删 tail 前节点并从 HashMap 移除对应 key};

\end{tikzpicture}
\end{document}
